\documentclass[11pt,a4paper]{article}
\usepackage[T1]{fontenc}
\usepackage[utf8]{inputenc}
\usepackage{lmodern}
\usepackage{amsmath,amssymb,amsthm}
\usepackage{booktabs}
\usepackage[margin=1.05in]{geometry}
\usepackage[dvipsnames]{xcolor}
\usepackage[colorlinks=true,linkcolor=blue!50!black,citecolor=blue!50!black,urlcolor=blue!50!black]{hyperref}
\hyphenpenalty=10000 \exhyphenpenalty=10000 \tolerance=4000 \emergencystretch=3em
\frenchspacing

\newtheorem{theoreme}{Théorème}
\newtheorem{proposition}[theoreme]{Proposition}
\newtheorem{lemme}[theoreme]{Lemme}
\theoremstyle{definition}
\newtheorem{conjecture}{Conjecture}
\theoremstyle{remark}
\newtheorem*{remarque}{Remarque}
\renewcommand{\proofname}{Démonstration}
\renewcommand{\refname}{Références}
\renewcommand{\abstractname}{Résumé}
\renewcommand{\tablename}{Table}

\newcommand{\R}{\mathbb{R}}
\newcommand{\dd}{\mathrm{d}}
\newcommand{\ps}[2]{\langle #1,\,#2\rangle}
\DeclareMathOperator{\re}{Re}
\DeclareMathOperator{\sh}{sh}
\DeclareMathOperator{\ch}{ch}

\title{La forme de Weil dans une fenêtre et la récurrence de Nicolas~:\\ ce qui se démontre, et ce qui manque}
\author{Antony Geremia\thanks{Chercheur indépendant, France. Note rédigée avec l'aide de Claude (Anthropic), dans le cadre du projet
\emph{La musique des nombres premiers}, \url{https://github.com/GEREMIAA/Musique-Des-Premiers}.}}
\date{Novembre 2026}

\begin{document}
\maketitle

\begin{abstract}
\textbf{Cette note ne démontre pas l'hypothèse de Riemann.} Elle rassemble, avec leurs démonstrations complètes, les énoncés
établis au cours d'une exploration numérique de deux reformulations de l'hypothèse~: la positivité de Weil, restreinte aux
fonctions portées par une fenêtre $[x^{-1/2},x^{1/2}]$, et le critère de Nicolas sur les primorielles. Pour la première, on
sépare la forme de Weil en $W=P+B$, où $P$ est le terme des pôles~; on montre que l'état fondamental de $B$ est strictement
positif (Perron--Frobenius), que les spectres de $W$ et de $B$ s'entrelacent, que la positivité de $W$ équivaut à deux
inégalités sur la résolvante de $B$, et que, \emph{si} l'hypothèse de Riemann est vraie, la plus petite énergie paire est
au moins $\gamma_1^2/R$ fois plus petite que la plus petite énergie impaire. Pour le second, on donne un critère exact de
décroissance de la marge de Nicolas d'une primorielle à la suivante~: elle décroît à chaque pas tant que $p\le 10^{19}$, et
croît une infinité de fois au-delà. Ces énoncés sont élémentaires ou classiques dans leur méthode, et probablement connus
des spécialistes sous d'autres formes. La dernière section dit précisément ce qui manque.
\end{abstract}

\section{La forme de Weil dans une fenêtre}

\subsection*{Notations}
Pour $\varphi\in L^2(\R)$ réelle à support compact, on note
\[
\widehat\varphi(z)=\int_{\R}\varphi(t)\,e^{izt}\,\dd t\quad(z\in\mathbb{C}),\qquad
g_\varphi(u)=\int_{\R}\varphi(t)\,\varphi(t+u)\,\dd t,\qquad (\tau_u\varphi)(t)=\varphi(t+u).
\]
On a $g_\varphi(0)=\|\varphi\|^2$ et $\|\varphi-\tau_u\varphi\|^2=2\,(g_\varphi(0)-g_\varphi(u))$. On note $\Lambda$ la
fonction de von Mangoldt, $\psi=\Gamma'/\Gamma$ la fonction digamma, $\gamma$ la constante d'Euler, et
$\rho=\tfrac12+iz_\rho$ les zéros non triviaux de $\zeta$, comptés avec multiplicité~; $\gamma_1=14{,}1347\ldots$ est
l'ordonnée du premier.

\subsection*{La formule explicite}
Pour $\varphi\in C_c^\infty(\R)$ réelle, la formule explicite de Guinand et Weil \cite{weil,bombieri} s'écrit
\begin{equation}\label{eq:explicite}
\sum_\rho \widehat\varphi(z_\rho)\,\widehat\varphi(-z_\rho)\;=\;W(\varphi),
\end{equation}
\begin{equation}\label{eq:W}
W(\varphi)=2\,\widehat\varphi(\tfrac i2)\,\widehat\varphi(-\tfrac i2)
+\frac1{2\pi}\int_{\R}|\widehat\varphi(r)|^2\Bigl(\re\psi\bigl(\tfrac14+\tfrac{ir}2\bigr)-\ln\pi\Bigr)\dd r
-2\sum_{n\ge2}\frac{\Lambda(n)}{\sqrt n}\,g_\varphi(\ln n).
\end{equation}
Si l'hypothèse de Riemann est vraie, les $z_\rho=\gamma$ sont réels et le membre de gauche vaut
$\sum_\gamma|\widehat\varphi(\gamma)|^2\ge0$. Le critère de Weil dit que la réciproque est vraie~: l'hypothèse de Riemann
équivaut à $W(\varphi)\ge0$ pour toute $\varphi$.

\subsection*{La décomposition $W=P+B$}
La représentation intégrale de Gauss, $\psi(z)=-\gamma+\int_0^\infty\frac{e^{-v}-e^{-zv}}{1-e^{-v}}\,\dd v$ pour $\re z>0$, donne, après
le changement de variable $v=2u$,
\begin{equation}\label{eq:levy}
\re\psi\bigl(\tfrac14+\tfrac{ir}2\bigr)=\psi(\tfrac14)+\int_0^\infty\bigl(1-\cos(ru)\bigr)\,\nu(u)\,\dd u,
\qquad \nu(u)=\frac{e^{u/2}}{\sh u}>0 .
\end{equation}
Comme $\frac1{2\pi}\int|\widehat\varphi(r)|^2(1-\cos ru)\,\dd r=g_\varphi(0)-g_\varphi(u)=\tfrac12\|\varphi-\tau_u\varphi\|^2$, on obtient
$W=P+B$ avec
\begin{align}
P(\varphi)&=2\,\widehat\varphi(\tfrac i2)\,\widehat\varphi(-\tfrac i2)
=2\int\varphi(t)e^{t/2}\dd t\int\varphi(t)e^{-t/2}\dd t,\label{eq:P}\\
B(\varphi)&=E(\varphi)-\kappa\,\|\varphi\|^2-2\sum_{n\ge2}\frac{\Lambda(n)}{\sqrt n}\,g_\varphi(\ln n),\qquad
E(\varphi)=\frac12\int_0^\infty\|\varphi-\tau_u\varphi\|^2\,\nu(u)\,\dd u,\label{eq:B}
\end{align}
où $\kappa=\ln\pi-\psi(\tfrac14)=\ln(8\pi)+\gamma+\tfrac\pi2\approx5{,}372$. Le terme $P$ vient des pôles de $\zeta$~; $E$ est la
forme de Dirichlet d'un processus de Lévy symétrique, de mesure de Lévy $\nu(u)\,\dd u$.

\subsection*{La fenêtre}
On fixe $x>1$, $a=\tfrac12\ln x$ et $I=\,]-a,a[$. Une fonction de $L^2(I)$, prolongée par $0$, est une «~onde vivant entre
$x^{-1/2}$ et $x^{1/2}$~» (variable $t=\ln$ de la position). Pour une telle $\varphi$, $g_\varphi(u)=0$ dès que $|u|\ge\ln x$~:
dans \eqref{eq:B}, seules les puissances de nombres premiers $n<x$ interviennent. On note
$\mathcal D=\{\varphi\in L^2(I): E(\varphi)<\infty\}$ le domaine de la forme. Le symbole
$\eta(r)=\re\psi(\frac14+\frac{ir}2)-\psi(\frac14)$ de $E$ tend vers $+\infty$ comme $\ln|r|$~; comme $I$ est borné,
l'injection de $\mathcal D$ dans $L^2(I)$ est compacte. La forme $B$ est fermée et minorée sur $\mathcal D$, et $P$ est
bornée~: les opérateurs autoadjoints associés à $B$ et à $W$, encore notés $B$ et $W$, sont à résolvante compacte. On note
$\lambda_0(\cdot)\le\lambda_1(\cdot)\le\cdots$ leurs valeurs propres. La réflexion $t\mapsto-t$ commute avec $B$ et $W$~;
on note $B_{\mathrm p},W_{\mathrm p}$ (resp. $B_{\mathrm i},W_{\mathrm i}$) leurs restrictions aux fonctions paires (resp.
impaires), et
\[
c(t)=\ch(t/2)\,\mathbf 1_I(t),\qquad s(t)=\sh(t/2)\,\mathbf 1_I(t).
\]

\section{Ce qui se démontre dans une fenêtre}

\begin{lemme}[pair et impair se séparent]\label{lem:parite}
Pour $\varphi$ paire, $P(\varphi)=2\ps{\varphi}{c}^2$~; pour $\varphi$ impaire, $P(\varphi)=-2\ps{\varphi}{s}^2$. Ainsi
\[
W_{\mathrm p}=B_{\mathrm p}+2\,|c\rangle\langle c|,\qquad W_{\mathrm i}=B_{\mathrm i}-2\,|s\rangle\langle s| .
\]
\end{lemme}
\begin{proof}
Si $\varphi$ est paire, $\int\varphi e^{\pm t/2}=\int\varphi\ch(t/2)$~; si elle est impaire, $\int\varphi e^{\pm t/2}=\pm\int\varphi\sh(t/2)$.
On conclut par \eqref{eq:P}.
\end{proof}

Les pôles relèvent donc les ondes paires et abaissent les ondes impaires.

\begin{proposition}[Perron--Frobenius sans les pôles]\label{prop:PF}
La plus petite valeur propre $\lambda_0(B)$ est simple, et sa fonction propre $\varphi_0$ peut être choisie strictement
positive presque partout sur $I$. En particulier $\varphi_0$ est paire.
\end{proposition}
\begin{proof}
Pour $\varphi\in\mathcal D$, on a $\bigl||\varphi(t)|-|\varphi(t+u)|\bigr|\le|\varphi(t)-\varphi(t+u)|$, donc
$|\varphi|\in\mathcal D$ et $E(|\varphi|)\le E(\varphi)$. De plus $g_{|\varphi|}(u)\ge g_\varphi(u)$ pour tout $u$, et
$\||\varphi|\|=\|\varphi\|$. Donc $B(|\varphi|)\le B(\varphi)$~: par le premier critère de Beurling et Deny
\cite[§XIII.12]{reedsimon}, le semi-groupe $e^{-tB}$ préserve la positivité.

Montrons qu'il est irréductible. Soit $S\subset I$ mesurable tel que $L^2(S)$ soit invariant par le semi-groupe. Alors la
multiplication par $\mathbf 1_S$ préserve $\mathcal D$ et la forme bilinéaire vérifie $B(\mathbf 1_S\varphi,\mathbf 1_{I\setminus S}\varphi)=0$
pour toute $\varphi\in\mathcal D$ \cite[chap.~2]{ouhabaz}. Prenons $\varphi=\mathbf 1_I$, qui est dans $\mathcal D$ car
$\|\mathbf 1_I-\tau_u\mathbf 1_I\|^2=2\min(u,2a)$ et $\int_0^1u\,\nu(u)\,\dd u<\infty$. Posons $f=\mathbf 1_S$,
$h=\mathbf 1_{I\setminus S}$. Comme $f$ et $h$ sont à supports disjoints, $\ps fh=0$ et
\[
B(f,h)=-\frac12\int_0^\infty\bigl(\ps{f}{\tau_uh}+\ps{\tau_uf}{h}\bigr)\nu(u)\,\dd u
-\sum_{n<x}\frac{\Lambda(n)}{\sqrt n}\bigl(\ps{f}{\tau_{\ln n}h}+\ps{h}{\tau_{\ln n}f}\bigr).
\]
Tous les termes sont négatifs ou nuls, et $\nu>0$. Si $B(f,h)=0$, alors
$\int_0^{2a}\!\int\bigl(f(t)h(t+u)+f(t+u)h(t)\bigr)\dd t\,\dd u=\iint_{I\times I}f(t)h(t')\,\dd t\,\dd t'=|S|\cdot|I\setminus S|=0$.
Donc $S$ ou son complémentaire est négligeable~: le semi-groupe est irréductible. Un semi-groupe préservant la positivité et
irréductible améliore la positivité, et sa plus petite valeur propre, qui existe ici puisque la résolvante est compacte, est
simple avec une fonction propre strictement positive \cite[§XIII.12]{reedsimon}. Enfin, $\varphi_0(-t)$ est
aussi fonction propre pour $\lambda_0(B)$, donc lui est proportionnelle, et le facteur vaut $1$ puisque $\varphi_0>0$.
\end{proof}

\begin{proposition}[entrelacement et règles de somme]\label{prop:entrelacement}
Notons $\mu_{\mathrm p}=\lambda_0(W_{\mathrm p})$ et $\mu_{\mathrm i}=\lambda_0(W_{\mathrm i})$ les plus petites énergies paire et impaire, et,
pour $\mu$ hors du spectre de $B_{\mathrm p}$, $G(\mu)=1+2\,\ps{c}{(B_{\mathrm p}-\mu)^{-1}c}$.
\begin{enumerate}
\item[(a)] $\lambda_0(B)<\mu_{\mathrm p}\le\lambda_1(B_{\mathrm p})$, et $\mu_{\mathrm i}\le\lambda_0(B_{\mathrm i})$.
\item[(b)] Si $\mu_{\mathrm p}<\lambda_1(B_{\mathrm p})$, alors $\mu_{\mathrm p}$ est une valeur propre simple de $W_{\mathrm p}$~; c'est l'unique zéro de $G$
dans $]\lambda_0(B),\lambda_1(B_{\mathrm p})[$, et sa fonction propre est $(B_{\mathrm p}-\mu_{\mathrm p})^{-1}c$. Le cas
$\mu_{\mathrm p}=\lambda_1(B_{\mathrm p})$ ne se produit que si $c$ est orthogonale à tout l'espace propre de $\lambda_1(B_{\mathrm p})$.
\item[(c)] Supposons $0$ hors du spectre de $B_{\mathrm p}$. Alors $W_{\mathrm p}\ge0$ si et seulement si $B_{\mathrm p}\ge 0$, ou bien
$B_{\mathrm p}$ a exactement une valeur propre négative et
\begin{equation}\label{eq:regle-paire}
\ps{c}{B_{\mathrm p}^{-1}c}\le-\tfrac12 .
\end{equation}
\item[(d)] Supposons $B_{\mathrm i}>0$. Alors $W_{\mathrm i}\ge0$ si et seulement si
\begin{equation}\label{eq:regle-impaire}
\ps{s}{B_{\mathrm i}^{-1}s}\le\tfrac12 .
\end{equation}
\end{enumerate}
\end{proposition}
\begin{proof}
Notons $\lambda_0<\lambda_1\le\lambda_2\le\cdots$ les valeurs propres de $B_{\mathrm p}$ ($\lambda_0=\lambda_0(B)$ est simple par la
proposition~\ref{prop:PF}) et $\varphi_0,\varphi_1,\ldots$ une base orthonormée de fonctions propres. Alors
$G(\mu)=1+2\sum_j\ps{c}{\varphi_j}^2/(\lambda_j-\mu)$, et $\ps c{\varphi_0}>0$ car $c>0$ et $\varphi_0>0$.

(a) Comme $W_{\mathrm p}\ge B_{\mathrm p}$, $\mu_{\mathrm p}\ge\lambda_0$~; l'égalité exigerait une fonction propre de $B_{\mathrm p}$ pour $\lambda_0$
orthogonale à $c$, ce qui est exclu. Dans le plan engendré par $\varphi_0$ et $\varphi_1$, il existe un vecteur unitaire $v$
orthogonal à $c$~; alors $W_{\mathrm p}(v)=B_{\mathrm p}(v)\le\lambda_1$, donc $\mu_{\mathrm p}\le\lambda_1$. Enfin $W_{\mathrm i}\le B_{\mathrm i}$.

(b) Soit $\mu\in\,]\lambda_0,\lambda_1[$. L'équation $W_{\mathrm p}\phi=\mu\phi$ s'écrit $(B_{\mathrm p}-\mu)\phi=-2\ps c\phi\,c$. Si $\ps c\phi=0$, $\phi$ serait
vecteur propre de $B_{\mathrm p}$ pour $\mu$~: impossible. Donc $\phi$ est proportionnel à $(B_{\mathrm p}-\mu)^{-1}c$, et en prenant
le produit scalaire avec $c$, $G(\mu)=0$. Réciproquement tout zéro de $G$ donne une valeur propre. Sur $]\lambda_0,\lambda_1[$,
$G'(\mu)=2\sum_j\ps c{\varphi_j}^2/(\lambda_j-\mu)^2>0$ et $G(\mu)\to-\infty$ quand $\mu\downarrow\lambda_0$~: $G$ y a au plus un zéro.
Si $\mu_{\mathrm p}<\lambda_1$, c'est ce zéro, et l'espace propre est de dimension $1$. Si $c$ n'est pas orthogonale à l'espace propre de
$\lambda_1$, alors $G(\mu)\to+\infty$ quand $\mu\uparrow\lambda_1$, $G$ s'annule dans l'intervalle, et $\mu_{\mathrm p}<\lambda_1$.

(c) Si $B_{\mathrm p}$ a au moins deux valeurs propres négatives (avec multiplicité), un plan où $B_{\mathrm p}<0$ contient un vecteur non nul
orthogonal à $c$, sur lequel $W_{\mathrm p}=B_{\mathrm p}<0$. Si $B_{\mathrm p}$ en a exactement une, alors $\lambda_0<0<\lambda_1$, et d'après (a) et (b)~:
$W_{\mathrm p}\ge0$ si et seulement si $\mu_{\mathrm p}\ge0$, si et seulement si $G$ n'a pas de zéro dans $]\lambda_0,0[$, c'est-à-dire, $G$
étant croissante, si et seulement si $G(0)\le0$. C'est \eqref{eq:regle-paire}.

(d) $W_{\mathrm i}=B_{\mathrm i}^{1/2}\bigl(1-2\,|w\rangle\langle w|\bigr)B_{\mathrm i}^{1/2}$ avec $w=B_{\mathrm i}^{-1/2}s$, et $1-2|w\rangle\langle w|\ge0$ si et
seulement si $2\|w\|^2\le1$.
\end{proof}

\begin{remarque}
D'après le critère de Weil, l'hypothèse de Riemann équivaut donc à ceci~: pour toute fenêtre, $B_{\mathrm p}$ a au plus une
valeur propre négative, $B_{\mathrm i}\ge0$, et les deux «~règles de somme~» \eqref{eq:regle-paire} et \eqref{eq:regle-impaire}
sont vraies (aux cas dégénérés près). Numériquement (table~\ref{tab:sommes}), elles le sont avec une marge qui fond~:
$\ps c{B_{\mathrm p}^{-1}c}=-\tfrac12-\varepsilon$ avec $\varepsilon\approx2\cdot10^{-30}$ pour $x=8$ et $4\cdot10^{-58}$ pour $x=64$.
Le terme des pôles compense exactement la seule direction négative de $B$.
\end{remarque}

\begin{proposition}[sous l'hypothèse de Riemann, une onde impaire donne une onde paire plus silencieuse]\label{prop:primitive}
Supposons l'hypothèse de Riemann. Soit $\phi\in C_c^\infty(I)$ impaire, et $\Phi(t)=\int_{-a}^t\phi$ sa primitive. Alors
$\Phi\in C_c^\infty(I)$ est paire et
\[
W(\Phi)\le\frac{W(\phi)}{\gamma_1^{2}}.
\]
Par conséquent $\mu_{\mathrm p}\le \dfrac{R(\phi)}{\gamma_1^2}\cdot\dfrac{W(\phi)}{\|\phi\|^2}$, où $R(\phi)=\|\phi\|^2/\|\Phi\|^2$.
\end{proposition}
\begin{proof}
Comme $\phi$ est impaire, $\int_I\phi=0$, donc $\Phi$ est nulle hors de $I$, lisse et paire. De $\phi=\Phi'$ on tire
$\widehat\phi(r)=-ir\,\widehat\Phi(r)$. Sous l'hypothèse de Riemann, \eqref{eq:explicite} donne
$W(\Phi)=\sum_\gamma|\widehat\Phi(\gamma)|^2=\sum_\gamma|\widehat\phi(\gamma)|^2/\gamma^2\le\gamma_1^{-2}\sum_\gamma|\widehat\phi(\gamma)|^2
=W(\phi)/\gamma_1^2$, puisque $|\gamma|\ge\gamma_1$ pour tout zéro. Enfin $\mu_{\mathrm p}\le W(\Phi)/\|\Phi\|^2$.
\end{proof}

Comme $\gamma_1^2\approx199{,}8$, l'onde la plus basse est paire dès qu'une onde impaire presque optimale a un rapport
$R<\gamma_1^2$~; dans nos calculs, $R$ vaut de $6{,}7$ à $16{,}3$ (table~\ref{tab:sommes}). Cette explication suppose
l'hypothèse~: elle n'en est pas une démonstration.

\subsection*{Le jeu de Weil et la loi de la marge}

\begin{theoreme}[le jeu de Weil n'a qu'un gagnant]\label{th:jeu}
Supposons l'hypothèse de Riemann. Soient $c_n$ des réels, égaux à $1$ sauf pour un nombre fini de puissances de nombres
premiers, et non tous égaux à $1$. Alors la forme $W_c$ obtenue en remplaçant $\Lambda(n)$ par $c_n\Lambda(n)$ dans
\eqref{eq:W} prend des valeurs strictement négatives.
\end{theoreme}
\begin{proof}
Par \eqref{eq:explicite} et l'identité $g_\varphi(u)=\frac1{2\pi}\int|\widehat\varphi(r)|^2\cos(ru)\,\dd r$,
\[
W_c(\varphi)=\sum_\gamma|\widehat\varphi(\gamma)|^2-\int_{\R}|\widehat\varphi(r)|^2F(r)\,\dd r,\qquad
F(r)=\frac1\pi\sum_n(c_n-1)\frac{\Lambda(n)}{\sqrt n}\cos(r\ln n).
\]
$F$ est une somme finie, non nulle, de cosinus de fréquences $\ln n>0$ distinctes~: elle est de moyenne nulle et non
identiquement nulle, donc strictement positive sur un intervalle ouvert. Les zéros étant isolés, on y choisit $r_0$ et
$\delta>0$ tels que $|r_0\mp\gamma|\ge\delta$ pour tout zéro. Soit $\chi\in C_c^\infty(\R)$ non nulle et
$\varphi_L(t)=L^{-1/2}\chi(t/L)\cos(r_0t)$. Alors
$\widehat\varphi_L(r)=\tfrac{\sqrt L}2\bigl(\widehat\chi(L(r-r_0))+\widehat\chi(L(r+r_0))\bigr)$. Quand $L\to\infty$,
$\int|\widehat\varphi_L|^2F\to\pi\|\chi\|^2F(r_0)>0$, tandis que, $\widehat\chi$ étant à décroissance rapide,
$\sum_\gamma|\widehat\varphi_L(\gamma)|^2\le C\,L^{-1}\sum_\gamma\bigl(|\gamma-r_0|^{-2}+|\gamma+r_0|^{-2}\bigr)\to0$.
Donc $W_c(\varphi_L)<0$ pour $L$ assez grand.
\end{proof}

\begin{lemme}[la loi de la marge]\label{lem:marge}
Dans une fenêtre, soit $\lambda$ la plus petite valeur propre de $W$ et $\varphi_1$ une fonction propre normalisée.
Pour une puissance de nombre premier $n<x$, posons $s_n=-2\Lambda(n)n^{-1/2}g_{\varphi_1}(\ln n)$. Si $s_n<0$, on ne peut pas
multiplier le poids de $n$ par plus de $1+\lambda/|s_n|$ sans que la forme devienne négative. Si $\varphi_1\ge0$, alors
$s_n\le0$ pour tout $n$. Si de plus $\varphi_1$ est paire et log-concave, $g_{\varphi_1}$ est décroissante sur $[0,\infty[$, et
la borne $\lambda/|s_p|$ est croissante le long des nombres premiers $p\ge7$.
\end{lemme}
\begin{proof}
Avec le poids $1+t$ pour $n$, l'énergie de $\varphi_1$ vaut $\lambda+t\,s_n$. Si $\varphi_1\ge0$, $g_{\varphi_1}\ge0$. Si
$\varphi_1$ est paire et log-concave, son autocorrélation est paire et log-concave (Prékopa--Leindler), donc décroissante sur
$[0,\infty[$~; et $2\ln p/\sqrt p$ décroît le long des nombres premiers à partir de $7$.
\end{proof}

\subsection*{Vérifications numériques}
Les tables \ref{tab:niveaux} et \ref{tab:sommes} donnent les quantités des propositions \ref{prop:entrelacement} et
\ref{prop:primitive}, calculées par une méthode de Galerkin avec $20$ ondes de chaque parité (cosinus et sinus nuls au bord de
la fenêtre), en arithmétique décimale de $120$ à $200$ chiffres (programme \texttt{structure\_de\_weil.py}). Les énergies
calculées majorent celles de la forme complète (principe de Rayleigh--Ritz)~; elles illustrent les énoncés et ne les démontrent pas. Le programme \texttt{controles\_de\_la\_note.py} contrôle par ailleurs, sur des ondes tests, la formule \eqref{eq:explicite}--\eqref{eq:W}, la représentation \eqref{eq:levy} et la décomposition $W=P+B$.
Dans toutes les fenêtres, l'état fondamental calculé de $B$ est strictement positif, $B_{\mathrm p}$ a une seule valeur
propre négative, $B_{\mathrm i}>0$, et $\lambda_0(B)<\mu_{\mathrm p}<\lambda_1(B_{\mathrm p})$, $\mu_{\mathrm i}<\lambda_0(B_{\mathrm i})$,
$\mu_{\mathrm p}<\mu_{\mathrm i}$.

\begin{table}[ht]\centering\small
\begin{tabular}{rrrrrr}\toprule
$x$ & $\lambda_0(B)$ & $\mu_{\mathrm p}$ & $\lambda_1(B_{\mathrm p})$ & $\mu_{\mathrm i}$ & $\lambda_0(B_{\mathrm i})$\\\midrule
$2{,}5$ & $-1{,}80$ & $1{,}2\cdot 10^{-5}$ & $5{,}9\cdot 10^{-2}$ & $1{,}9\cdot 10^{-3}$ & $2{,}0\cdot 10^{-2}$ \\
$4$ & $-2{,}81$ & $9{,}1\cdot 10^{-13}$ & $1{,}2\cdot 10^{-7}$ & $5{,}7\cdot 10^{-10}$ & $1{,}4\cdot 10^{-5}$ \\
$8$ & $-4{,}46$ & $5{,}8\cdot 10^{-30}$ & $4{,}2\cdot 10^{-24}$ & $4{,}9\cdot 10^{-27}$ & $8{,}9\cdot 10^{-22}$ \\
$16$ & $-6{,}42$ & $3{,}4\cdot 10^{-42}$ & $2{,}7\cdot 10^{-36}$ & $1{,}2\cdot 10^{-39}$ & $1{,}7\cdot 10^{-34}$ \\
$32$ & $-8{,}78$ & $1{,}2\cdot 10^{-50}$ & $7{,}3\cdot 10^{-45}$ & $1{,}5\cdot 10^{-48}$ & $3{,}3\cdot 10^{-43}$ \\
$64$ & $-11{,}79$ & $1{,}7\cdot 10^{-57}$ & $1{,}7\cdot 10^{-51}$ & $2{,}8\cdot 10^{-55}$ & $3{,}3\cdot 10^{-50}$ \\\bottomrule
\end{tabular}
\caption{Entrelacement des énergies de $W$ et de $B$ ($20$ ondes de chaque parité).}\label{tab:niveaux}
\end{table}

\begin{table}[ht]\centering\small
\begin{tabular}{rrrrrr}\toprule
$x$ & $-1-2\ps c{B_{\mathrm p}^{-1}c}$ & $1-2\ps s{B_{\mathrm i}^{-1}s}$ & $\mu_{\mathrm i}/\mu_{\mathrm p}$ & $R$ & $\gamma_1^2/R$\\\midrule
$2{,}5$ & $9{,}6\cdot 10^{-6}$ & $9{,}5\cdot 10^{-2}$ & $163$ & $16{,}3$ & $12{,}3$ \\
$4$ & $6{,}6\cdot 10^{-13}$ & $2{,}1\cdot 10^{-8}$ & $626$ & $12{,}9$ & $15{,}5$ \\
$8$ & $3{,}9\cdot 10^{-30}$ & $1{,}5\cdot 10^{-25}$ & $844$ & $11{,}2$ & $17{,}8$ \\
$16$ & $2{,}1\cdot 10^{-42}$ & $2{,}8\cdot 10^{-38}$ & $339$ & $9{,}6$ & $20{,}7$ \\
$32$ & $6{,}9\cdot 10^{-51}$ & $2{,}7\cdot 10^{-47}$ & $119$ & $8{,}1$ & $24{,}8$ \\
$64$ & $8{,}5\cdot 10^{-58}$ & $3{,}9\cdot 10^{-54}$ & $170$ & $6{,}7$ & $29{,}8$ \\\bottomrule
\end{tabular}
\caption{Les marges des deux règles de somme \eqref{eq:regle-paire} et \eqref{eq:regle-impaire} (colonnes 2 et 3, positives),
et le rapport des plus petites énergies impaire et paire, comparé à la borne $\gamma_1^2/R$ de la proposition~\ref{prop:primitive},
où $R$ est calculé sur l'onde impaire la plus basse.}\label{tab:sommes}
\end{table}

\section{La récurrence de Nicolas}

Soient $p_1=2<p_2=3<\cdots$ les nombres premiers, $\theta(x)=\sum_{p\le x}\ln p$, $\theta_k=\theta(p_k)$, et, pour $k\ge2$,
\[
F_k=\sum_{j\le k}-\ln\Bigl(1-\frac1{p_j}\Bigr)-\gamma-\ln\ln\theta_k .
\]
Avec $N_k=p_1\cdots p_k$ et $\varphi$ l'indicatrice d'Euler, $F_k=\ln\bigl(N_k/\varphi(N_k)\bigr)-\ln\bigl(e^\gamma\ln\ln N_k\bigr)$.
Nicolas \cite{nicolas} a démontré que l'hypothèse de Riemann équivaut à $F_k>0$ pour tout $k$~; et que, si elle est fausse,
$F_k$ change de signe une infinité de fois.

\begin{proposition}[la marge est une intégrale sur l'avenir]\label{prop:avenir}
Pour $k\ge2$,
\[
F_k=\int_{p_k}^\infty\bigl(t-\theta(t)\bigr)\,\frac{1+\ln t}{t^2\ln^2t}\,\dd t\;-\;H_k\;+\;\varepsilon_k,
\]
où $H_k=\sum_{q>p_k}\bigl(-\ln(1-\tfrac1q)-\tfrac1q\bigr)\in\,]0,\tfrac1{2p_k}[$ et
$\varepsilon_k=\ln\ln p_k-\ln\ln\theta_k-\dfrac{p_k-\theta_k}{p_k\ln p_k}\ge0$.
\end{proposition}
\begin{proof}
Posons $w(t)=(1+\ln t)/(t^2\ln^2t)=-\frac{\dd}{\dd t}\frac1{t\ln t}$. Par sommation d'Abel,
$\sum_{q\le x}\frac1q=\frac{\theta(x)}{x\ln x}+\int_2^x\theta(t)w(t)\,\dd t$, et $\int_2^xt\,w(t)\,\dd t=\ln\ln x-\frac1{\ln x}-\ln\ln2+\frac1{\ln2}$.
Par le théorème des nombres premiers, $\int_2^\infty(\theta(t)-t)w(t)\,\dd t$ converge, d'où
$\sum_{q\le x}\frac1q=\ln\ln x+M+\frac{\theta(x)-x}{x\ln x}+\int_x^\infty(t-\theta(t))w(t)\,\dd t$, où $M$ est la constante de
Meissel--Mertens. Le théorème de Mertens donne $\gamma=M+\sum_q\bigl(-\ln(1-\frac1q)-\frac1q\bigr)$. En $x=p_k$, on obtient la
formule annoncée. L'encadrement de $H_k$ vient de $-\ln(1-\frac1q)-\frac1q<\frac1{2q(q-1)}$, et $\varepsilon_k\ge0$ de la convexité de
$D\mapsto-\ln\ln(p_k-D)$, dont la dérivée en $0$ vaut $1/(p_k\ln p_k)$.
\end{proof}

La valeur de $F_k$ se calcule avec les $k$ premiers nombres premiers, mais elle est égale à une moyenne pondérée du déficit
$t-\theta(t)$ \emph{au-delà} de $p_k$~: c'est la constante $\gamma$ qui contient cet avenir.

\begin{lemme}[critère exact de décroissance]\label{lem:critere}
Soit $k\ge2$, $p=p_{k+1}$ et $\theta^*(p)$ l'unique solution $\theta>1$ de $(p-1)\ln\bigl(1+\frac{\ln p}{\theta}\bigr)=\ln\theta$. Alors
\[
F_{k+1}<F_k\iff\theta_k<\theta^*(p_{k+1}),\qquad\text{et}\qquad p-\tfrac12\ln p-\tfrac34<\theta^*(p)<p .
\]
\end{lemme}
\begin{proof}
On a $F_{k+1}-F_k=-\ln(1-\frac1p)-\ln\frac{\ln(\theta_k+\ln p)}{\ln\theta_k}$. Comme $\ln\theta_k>0$, ceci est négatif si et seulement si
$\Phi(\theta_k)>0$, où $\Phi(\theta)=(p-1)\ln(1+\frac{\ln p}\theta)-\ln\theta$ est strictement décroissante. D'une part
$\Phi(p)<(p-1)\frac{\ln p}p-\ln p<0$. D'autre part, soit $L=\ln p$, $\delta_0=\frac L2+\frac34$ et $\theta_0=p-\delta_0>0$. De
$\ln(1+y)\ge\frac{2y}{2+y}$ et $\ln\theta_0\le L-\frac{\delta_0}p$, il suffit, pour avoir $\Phi(\theta_0)>0$, que
$2L(p-1)>(2\theta_0+L)(L-\frac{\delta_0}p)=(2p-\frac32)(L-\frac{\delta_0}p)$~; la différence des deux membres vaut
$\frac L2+\frac32-\frac{3\delta_0}{2p}>0$.
\end{proof}

\begin{theoreme}[la marge de Nicolas décroît longtemps, mais pas toujours]\label{th:nicolas}
\leavevmode
\begin{enumerate}
\item[(a)] $F_{k+1}<F_k$ pour tout $k\ge2$ tel que $p_{k+1}\le10^{19}$.
\item[(b)] $F_{k+1}>F_k$ pour une infinité de $k$.
\item[(c)] Le plus petit $k$ tel que $F_{k+1}>F_k$ vérifie $10^{19}<p_{k+1}<e^{727{,}951332668}<1{,}4\cdot10^{316}$.
\end{enumerate}
\end{theoreme}
\begin{proof}
(a) Büthe \cite[th.~2]{buthe} a démontré $x-\theta(x)>0{,}05\sqrt x$ pour $1\le x\le10^{19}$. En faisant tendre $x$ vers $p_{k+1}$ par
valeurs inférieures, $p_{k+1}-\theta_k\ge0{,}05\sqrt{p_{k+1}}$, et $0{,}05\sqrt p\ge\frac12\ln p+\frac34$ dès que $p\ge11\,835$. Le
lemme~\ref{lem:critere} conclut pour $p_{k+1}\ge11\,835$. Les $1\,436$ pas restants ($5\le p_{k+1}<12\,000$) se vérifient par un calcul
direct en $60$ chiffres (programme \texttt{nicolas\_petits\_premiers.py})~: le moins net est $F_{k+1}-F_k=-7{,}6\cdot10^{-7}$, en $p_{k+1}=10\,343$.

(b) Littlewood \cite{littlewood} a démontré $\psi(x)-x=\Omega_\pm(\sqrt x\,\ln\ln\ln x)$~; comme $\psi(x)-\theta(x)=O(\sqrt x)$, il
existe des $x$ arbitrairement grands avec $\theta(x)>x+\ln x$. Soit $p_m$ le plus grand nombre premier $\le x$. Alors
$\theta_{m-1}=\theta(x)-\ln p_m>x\ge p_m>\theta^*(p_m)$, donc $F_m>F_{m-1}$ par le lemme~\ref{lem:critere}.

(c) Platt et Trudgian \cite{platttrudgian} ont montré qu'il existe $x\in[e^{727{,}951332642},e^{727{,}951332668}]$ avec
$\theta(x)-x>10^{152}$. L'argument de (b) s'applique à ce $x$~; la borne inférieure vient de (a).
\end{proof}

\begin{remarque}
Le crible du projet (programme \texttt{nicolas\_premier\_apres\_premier.c}) a suivi $F_k$ pour les $4\,118\,054\,813$ nombres premiers
jusqu'à $10^{11}$~: aucune remontée, conformément à (a). L'énoncé «~la marge baisse à chaque pas~» est donc vrai pour les
$234\,057\,667\,276\,344\,607$ premiers nombres premiers (ceux qui précèdent $10^{19}$), et faux une infinité de fois~: c'est un
exemple de plus de ce qu'une vérification finie ne dit pas.
\end{remarque}

\section{Ce qui manque}

Le plan le plus avancé pour démontrer l'hypothèse de Riemann à partir de la forme de Weil d'une fenêtre est celui de Connes,
Consani et Moscovici \cite{ccm}~: si la plus petite valeur propre de leur forme restreinte est simple et de fonction propre paire,
la transformée de Fourier de cette fonction propre n'a que des zéros réels~; si de plus ces transformées convergent vers la
fonction $\Xi$ de Riemann quand la fenêtre grandit, le théorème de Hurwitz donne l'hypothèse. Les énoncés de la section~2 ne
franchissent aucune de ces deux étapes. Ils les ramènent à ceci.

\begin{conjecture}[simplicité]
Dans toute fenêtre, $\mu_{\mathrm p}<\lambda_1(B_{\mathrm p})$~; autrement dit, $G$ s'annule dans
$]\lambda_0(B),\lambda_1(B_{\mathrm p})[$, ce qui est automatique dès que $c$ n'est pas orthogonale à l'espace propre de
$\lambda_1(B_{\mathrm p})$.
\end{conjecture}
\begin{conjecture}[parité]
Dans toute fenêtre, $\mu_{\mathrm p}<\mu_{\mathrm i}$.
\end{conjecture}
\begin{conjecture}[convergence]
Quand $x\to\infty$, la fonction propre de $\mu_{\mathrm p}$, normalisée, converge vers la fonction $\Phi$ de Riemann, dont $\Xi$ est la
transformée de Fourier, assez fortement pour que les transformées de Fourier convergent uniformément sur les compacts de
$\mathbb{C}$.
\end{conjecture}

La proposition~\ref{prop:primitive} montre que la conjecture de parité découle de l'hypothèse de Riemann (quand $R<\gamma_1^2$)~:
elle est donc au moins compatible avec elle, mais nous ne savons pas la démontrer sans elle. La difficulté commune se lit dans la
table~\ref{tab:niveaux}~: toutes les quantités à comparer sont de l'ordre de $10^{-30}$ à $10^{-57}$, alors que $B$ a des
valeurs propres d'ordre $1$. Aucune fonction d'essai grossière ne peut servir~; il faut des identités exactes. La conjecture de
convergence contient, à elle seule, l'essentiel de l'hypothèse de Riemann.

\begin{thebibliography}{9}
\bibitem{bombieri} E. Bombieri, \emph{Remarks on Weil's quadratic functional in the theory of prime numbers I}, Rend. Mat. Acc. Lincei 11 (2000), 183--233.
\bibitem{buthe} J. Büthe, \emph{An analytic method for bounding $\psi(x)$}, Math. Comp. 87 (2018), 1991--2009, \href{https://arxiv.org/abs/1511.02032}{arXiv:1511.02032}.
\bibitem{ccm} A. Connes, C. Consani, H. Moscovici, \emph{Zeta spectral triples}, \href{https://arxiv.org/abs/2511.22755}{arXiv:2511.22755} (2025).
\bibitem{littlewood} J. E. Littlewood, \emph{Sur la distribution des nombres premiers}, C. R. Acad. Sci. Paris 158 (1914), 1869--1872.
\bibitem{nicolas} J.-L. Nicolas, \emph{Petites valeurs de la fonction d'Euler}, J. Number Theory 17 (1983), 375--388.
\bibitem{ouhabaz} E. M. Ouhabaz, \emph{Analysis of Heat Equations on Domains}, Princeton University Press, 2005.
\bibitem{platttrudgian} D. J. Platt, T. S. Trudgian, \emph{On the first sign change of $\theta(x)-x$}, Math. Comp. 85 (2016), 1539--1547, \href{https://arxiv.org/abs/1407.1914}{arXiv:1407.1914}.
\bibitem{reedsimon} M. Reed, B. Simon, \emph{Methods of Modern Mathematical Physics IV: Analysis of Operators}, Academic Press, 1978.
\bibitem{weil} A. Weil, \emph{Sur les «~formules explicites~» de la théorie des nombres premiers}, Comm. Sém. Math. Univ. Lund (1952), 252--265.
\end{thebibliography}
\end{document}
